\chapter{The six Haagerup lines}
\label{ch:haagerup}

We now turn the abstract language of fusion into a concrete algebra.
The Haagerup example is small enough that its entire fusion ring can
be understood by hand.  Its associator already contains information
that no group multiplication table can express.  We will derive its
quantum dimensions, all irreducible representations of its ordinary
fusion algebra, an exact invariant-state projector, and a useful
four-dimensional recoupling matrix.  These calculations also fix the
notation used throughout the construction.

\section{A six-object fusion ring}
\index{Haagerup fusion ring}

Let $\alpha$ be an invertible object of order three and $\rho$ a
noninvertible object.  The simple objects are
\[
 \one,\alpha,\alpha^2,\rho,\alpha\rho,\alpha^2\rho.
\]
Their fusion rules are generated by
\begin{equation}
 \alpha^3=\one,\qquad
 \rho\alpha=\alpha^{-1}\rho,\qquad
 \rho^2=\one+\rho+\alpha\rho+\alpha^2\rho.
 \label{eq:H3-generating-ring}
\end{equation}
Direct sums are written with plus signs in a fusion ring.  Thus the
last equality means an isomorphism with a four-term direct sum.

Put $A_g=\alpha^g$ and $R_g=\alpha^g\rho$, with
$g\in\Z/3\Z$.  Moving invertible objects past $\rho$ gives the
complete table in four formulas:
\begin{align}
 A_g A_h&=A_{g+h},& A_gR_h&=R_{g+h},\notag\\
 R_gA_h&=R_{g-h},&
 R_gR_h&=A_{g-h}+R_0+R_1+R_2.
 \label{eq:H3-all-products}
\end{align}
For example,
\[
 R_1R_2=A_2+R_0+R_1+R_2,
 \qquad R_2R_1=A_1+R_0+R_1+R_2.
\]
The ring is not commutative.  This has an immediate categorical
consequence: a category with these fusion rules cannot admit a
braiding.  A braiding would give an isomorphism $a\otimes b\simeq
b\otimes a$ for every pair, making its Grothendieck ring commutative.
The braided category relevant to the CFT is the \emph{center} of the
Haagerup category, not the six-object category itself.

The duality operation follows by looking for the unit in a product:
\begin{equation}
 A_g^*=A_{-g},\qquad R_g^*=R_g.
 \label{eq:H3-duals}
\end{equation}
Every noninvertible simple is self-dual, even though the full ring is
noncommutative.

\begin{example}[An associativity check]
Using \eqref{eq:H3-all-products},
\begin{align*}
 (R_gR_h)R_k
 &=R_{g-h+k}+\sum_xA_{x-k}+3\sum_yR_y\\
 &=R_{g-h+k}+\sum_xA_x+3\sum_yR_y.
\end{align*}
Expanding $R_g(R_hR_k)$ gives the same expression.  This verifies a
representative family of associativity equations for the ring.
It does not verify the pentagon for an associator.  Ring associativity
counts channels, while the pentagon compares linear maps between
their bases.
\end{example}

\section{The number \texorpdfstring{$\sqrt{13}$}{sqrt(13)} is forced by fusion}

In a positive dimension function each invertible object has dimension
one, and multiplication by $\alpha$ shows that all three $R_g$ have
the same dimension $d$.  Applying the dimension function to
\eqref{eq:H3-generating-ring} gives
\[
 d^2=1+3d.
\]
Thus
\begin{equation}
 d=\frac{3+\sqrt{13}}2,\qquad
 d_-:=\frac{3-\sqrt{13}}2=-\frac1d.
 \label{eq:H3-dimension-roots}
\end{equation}
Only $d$ is positive.  We will frequently use the abbreviations
\begin{equation}
 s=1+d^2=2+3d=\frac{13+3\sqrt{13}}2,
 \qquad \mathcal D=3s.
 \label{eq:H3-s-D}
\end{equation}
The global dimension of the six-object category is
\begin{equation}
 \operatorname{Dim}(\HH_3)
 =\sum_x d_x^2=3+3d^2=3s=\mathcal D.
 \label{eq:H3-global-dimension}
\end{equation}
The symbol $\mathcal D$ will also be the total quantum dimension of
the twelve-object center.  The center's global dimension is
$\mathcal D^2$.  This is why distinguishing a global dimension from
its square root is essential.

The notation $d_-$ should not suggest a second positive quantum
dimension inside the same unitary category.  It is another algebraic
root and another scalar representation of the fusion ring.  A physical
unitary theory can contain states on which the closed-line operator
$Y_\rho$ has this negative
eigenvalue; an operator eigenvalue need not be a categorical dimension.

\section{One fusion ring, two unitary categories}
\index{Haagerup categories}

There are two unitary Haagerup categories customarily denoted
$\HH_2$ and $\HH_3$ with the six-object ring above.  Their associators
are inequivalent.  They are Morita equivalent, and consequently their
centers are braided equivalent, but these facts do not identify their
microscopic anyon chains.

The regular $\rho$-projector chain investigated later in the book uses
$\HH_3$. The $\HH_2$ presentation commonly used in algebraic
Haagerup--Izumi calculations has a different invariant coefficient,
compared below. The all-odd-order construction in
Chapter~\ref{ch:arithmetic} requires its own category and positivity
analysis. A change of label on a computer program cannot convert
one associator into the other.

\begin{constructioninput}[The categorical Haagerup input]
We use the established unitary category $\HH_3$ with fusion rule
\eqref{eq:H3-generating-ring} and its transparent associator
presentation.  ``Transparent'' here means a gauge in which admissible
associators involving an external invertible simple object are
identity maps.  It is a statement about an associator gauge, not a
claim that $\HH_3$ is braided or that a line lies in a M\"uger center.
The complete presentation satisfies the pentagon.  We display and
verify one important block below; that block alone is not a
replacement for the complete pentagon data.
\end{constructioninput}

\section{An explicit Haagerup recoupling matrix}

Consider three $\rho$ lines with total charge $\rho$.  Every summand
of $\rho^2$ is an admissible intermediate channel.  Use the order
\[
 (\one,\rho,\alpha\rho,\alpha^2\rho).
\]
Write $r=\sqrt{13}$ and define
\begin{equation}
 x=\frac{2-r}{3},\qquad
 y_- =\frac{5-r-\sqrt{6(1+r)}}{12},\qquad
 y_+ =\frac{5-r+\sqrt{6(1+r)}}{12}.
 \label{eq:H3-F-coefficients}
\end{equation}
In the convention of \eqref{eq:F-basis-convention}, the matrix is
\begin{equation}
 F^{\rho\rho\rho}_{\rho}
 =\begin{pmatrix}
 d^{-1}&d^{-1/2}&d^{-1/2}&d^{-1/2}\\
 d^{-1/2}&x&y_-&y_+\\
 d^{-1/2}&y_-&y_+&x\\
 d^{-1/2}&y_+&x&y_-
 \end{pmatrix}.
 \label{eq:H3-F-block}
\end{equation}
Changing the sign of the second radical exchanges the roles of
$\alpha\rho$ and $\alpha^2\rho$ in the corresponding full
presentation.  The displayed matrix is real and symmetric.

It is worthwhile to check its orthogonality entirely by hand.  First,
\begin{equation}
 x+y_-+y_+=\frac{3-r}{2}=-\frac1d.
 \label{eq:H3-F-sum}
\end{equation}
Next,
\[
 y_-y_+=\frac{(5-r)^2-6(1+r)}{144}
 =\frac{2-r}{9}=\frac{x}{3}.
\]
Consequently
\begin{align}
 xy_-+y_-y_++y_+x
 &=x\left(\frac{5-r}{6}+\frac13\right)
 =\frac{(2-r)(7-r)}{18}
 =-\frac1d.\label{eq:H3-F-pairs}
\end{align}
Squaring \eqref{eq:H3-F-sum} and using
\eqref{eq:H3-F-pairs} now gives
\begin{equation}
 x^2+y_-^2+y_+^2=d^{-2}+2d^{-1}=1-d^{-1}.
 \label{eq:H3-F-squares}
\end{equation}
The last equality is $d^{-2}+3d^{-1}=1$, which follows directly
from the dimension equation.

The first row of \eqref{eq:H3-F-block} has squared norm
$d^{-2}+3d^{-1}=1$.  Its inner product with any other row is
$d^{-1/2}(d^{-1}+x+y_-+y_+)=0$.  Every remaining row has squared
norm $d^{-1}+x^2+y_-^2+y_+^2=1$.  Finally, the inner product
of any two distinct remaining rows is
$d^{-1}+xy_-+y_-y_++y_+x=0$.  Hence
\begin{equation}
 (F^{\rho\rho\rho}_{\rho})^{\mathsf t}
 F^{\rho\rho\rho}_{\rho}=I.
 \label{eq:H3-F-orthogonal}
\end{equation}
This is an exact radical calculation; no decimal tolerance appears.

\begin{example}[Two explicit local interactions]
The right-vacuum projector in this left basis is
\[
 P^{(\one)}=
 \begin{pmatrix}
 d^{-2}&d^{-3/2}&d^{-3/2}&d^{-3/2}\\
 d^{-3/2}&d^{-1}&d^{-1}&d^{-1}\\
 d^{-3/2}&d^{-1}&d^{-1}&d^{-1}\\
 d^{-3/2}&d^{-1}&d^{-1}&d^{-1}
 \end{pmatrix}.
\]
The right-$\rho$ projector is the outer product of
\[
 v_\rho=(d^{-1/2},x,y_-,y_+)^{\mathsf t}
\]
with itself.  Both are rank one.  Their product vanishes by
\eqref{eq:H3-F-sum}.  Thus $-P^{(\rho)}$ and $-P^{(\one)}$
favor different fusion channels despite acting on the same four
local path configurations.  A periodic chain uses related blocks for
every possible pair of exterior path labels, not only this block.
\end{example}

The entry $x$ has an additional use.  It is exactly
$[F^{\rho\rho\rho}_{\rho}]_{\rho\rho}$, the gauge-invariant
entry discussed in Chapter~\ref{ch:fusion}.  For $\HH_2$ the same
entry is
\[
 x_{\HH_2}=\frac{7-\sqrt{13}}6,
 \qquad x_{\HH_3}=\frac{2-\sqrt{13}}3.
\]
Their difference is
\begin{equation}
 x_{\HH_2}-x_{\HH_3}
 =\frac{3+\sqrt{13}}6=\frac d3\ne0.
 \label{eq:H2-H3-invariant}
\end{equation}
They therefore cannot be related by a trivalent-basis gauge change
with this distinguished object identified.  Relabelling does not
remove the distinction: for every $R_g$, the entry
$[F^{R_gR_gR_g}_{R_g}]_{R_gR_g}$ is the same core entry
$[F^{\rho\rho\rho}_{\rho}]_{\rho\rho}$ within each category.
A fusion-rule permutation must take a noninvertible simple to a noninvertible
simple, so it preserves this comparison.  This concrete obstruction
explains why a result proved for an $\HH_2$ Hamiltonian cannot simply
be relabelled as a result for the $\HH_3$ chain.

\section{All representations of the ordinary fusion algebra}
\index{Grothendieck ring}

The Grothendieck ring $K_0(\HH_3)$ is the free abelian group on the
six simple-object classes, with addition induced by direct sum and
multiplication induced by fusion. Allowing complex coefficients gives
the six-dimensional associative algebra
$\mathscr A=K_0(\HH_3)\otimes_\Z\C$. It has two one-dimensional
representations:
\begin{equation}
 \alpha\mapsto1,\qquad \rho\mapsto d
 \quad\hbox{or}\quad
 \alpha\mapsto1,\qquad \rho\mapsto d_-.
 \label{eq:H3-scalar-reps}
\end{equation}
There is also a two-dimensional representation.  Put
$\omega=e^{2\pi i/3}$ and set
\begin{equation}
 A=\begin{pmatrix}\omega&0\\0&\omega^2\end{pmatrix},
 \qquad R=\begin{pmatrix}0&1\\1&0\end{pmatrix}.
 \label{eq:H3-doublet-rep}
\end{equation}
Then $A^3=I$, $RA=A^{-1}R$, and
\[
 I+R+AR+A^2R=I+(I+A+A^2)R=I=R^2.
\]
Thus the fusion rules hold.  The two eigenspaces of $A$ are exchanged
by $R$, so no one-dimensional subspace is invariant under both; this
representation is irreducible.

\begin{proposition}
There is an algebra isomorphism
\begin{equation}
 \mathscr A\simeq\C\oplus\C\oplus M_2(\C).
 \label{eq:H3-fusion-Wedderburn}
\end{equation}
The two scalar summands are the dimension and negative-root
characters, and the matrix summand is \eqref{eq:H3-doublet-rep}.
\end{proposition}
\begin{proof}
The three representations define a homomorphism into the displayed
algebra.  The spectral projectors of $A$ separate the two basis
directions in the matrix summand, and multiplication by $R$ gives
its two off-diagonal matrix units.  Thus that summand is generated.
The element
$E=(\one+\alpha+\alpha^2)/3$ vanishes in the matrix summand
and equals one in the two scalar summands.  The elements $E$ and
$E\rho$ separate those scalar summands because $d\ne d_-$.
The map is therefore surjective.  Its source and target both have
dimension $6=1+1+4$, so it is an isomorphism.
\end{proof}

This is the entire representation theory of the ordinary closed-ring
fusion algebra.  It has three irreducible representations, not twelve.
The twelve center sectors encountered later belong to the tube
algebra, which includes defect sectors and junction operations.

\section{An exact projector onto invariant states}
\index{invariant projector}

Suppose the line operators $Y_x$ on a Hilbert space obey the fusion
rules and adjoint relations $Y_x^\dagger=Y_{x^*}$. Define
\begin{equation}
 E=\frac{I+Y_\alpha+Y_{\alpha^2}}3.
 \label{eq:H3-neutral-projector}
\end{equation}
This is the ordinary $\Z_3$-neutral projector.  Since inversion
permutes the group elements, $E$ commutes with $Y_\rho$.  On its
image,
\[
 Y_\rho^2=I+3Y_\rho.
\]
The two possible eigenvalues are therefore $d$ and $d_-=-d^{-1}$.
Lagrange interpolation produces their projectors:
\begin{align}
 P_+&=E\frac{Y_\rho-d_-}{d-d_-}
 =E\frac{I+dY_\rho}{1+d^2},\label{eq:H3-Pplus}\\
 P_-&=E-P_+.
\end{align}
The equality of the two expressions follows after multiplying
numerator and denominator by $d$.

Expanding $E$ yields a particularly useful formula:
\begin{equation}
 P_+=\frac{
 I+Y_\alpha+Y_{\alpha^2}
 +d(Y_\rho+Y_{\alpha\rho}+Y_{\alpha^2\rho})}{3s}.
 \label{eq:H3-dimension-projector}
\end{equation}
It is the dimension-weighted average of all six line operators.
For a finite group every dimension equals one and the analogous
formula is the familiar group average.

\begin{proposition}
$P_+$ is the orthogonal projector onto states $v$ satisfying
\[
 Y_xv=d_xv\qquad\hbox{for all six simple lines }x.
\]
The projector $P_-$ selects states neutral under $\alpha$ but not
invariant under the full Haagerup symmetry.
\end{proposition}
\begin{proof}
On $\operatorname{im}E$, \eqref{eq:H3-Pplus} is one at the
eigenvalue $d$ and zero at $d_-$.  On $\ker E$ it vanishes.
Since $Y_\rho$ is self-adjoint and $E$ is an orthogonal projector,
this is an orthogonal spectral projector.  Its image has
$Y_\alpha=I$ and $Y_\rho=dI$, so
$Y_{\alpha^g\rho}=dI$ as well.  Conversely, a vector transforming
by all these dimensions lies in $\operatorname{im}E$ with
$Y_\rho=d$, hence is fixed by $P_+$.
\end{proof}

The charged central projector is $P_{\rm ch}=I-E$.  These three
projectors are pairwise orthogonal and sum to $I$.  In the regular
representation their ranks are $1,1,4$, respectively.  In a physical
Hilbert space each irreducible representation can occur many times,
at different energies.  A rank-one central block of the regular
algebra therefore does not imply a unique invariant physical state.

There is a further locality caveat with a simple explanation.
$P_+$ is built from lines wrapping the entire circle.  Sandwiching
a local operator $O$ as $P_+OP_+$ can turn it into a nonlocal
operator.  A projector onto invariant \emph{states} is not by itself
a recipe for every local invariant Hamiltonian density.

\section{Fusion paths and exact integer growth}

A periodic anyon chain with physical object $\rho$ has basis labels
$x_0,x_1,\ldots,x_L=x_0$ satisfying
$N_{x_j\rho}^{\ x_{j+1}}\ne0$.  For multiplicity-free fusion,
each admissible path contributes one basis vector.  In the order
$(A_0,A_1,A_2,R_0,R_1,R_2)$, right multiplication by $\rho$ has
adjacency matrix
\begin{equation}
 N_\rho^{\rm right}=
 \begin{pmatrix}0&I_3\\I_3&J_3\end{pmatrix},
 \label{eq:H3-path-adjacency}
\end{equation}
where $J_3$ has every entry equal to one.  Hence the number of
length-$L$ periodic paths is $\Tr((N_\rho^{\rm right})^L)$.

On the two-dimensional subspace of vectors constant in each block,
\eqref{eq:H3-path-adjacency} reduces to
$\left(\begin{smallmatrix}0&1\\1&3\end{smallmatrix}\right)$,
with eigenvalues $d,d_-$.  On the sum-zero subspace of each
three-dimensional block, $J_3$ vanishes, leaving eigenvalues $+1$
and $-1$, each twice.  Thus
\begin{equation}
 \dim\mathscr H_L=d^L+d_-^L+2(1+(-1)^L).
 \label{eq:H3-path-count}
\end{equation}
The irrational powers combine to give an integer.  Indeed, if
$C_L=d^L+d_-^L$, then
\begin{equation}
 C_0=2,\qquad C_1=3,\qquad C_L=3C_{L-1}+C_{L-2}.
 \label{eq:H3-path-recurrence}
\end{equation}
For $L=1,2,3,4,5,6$, the path counts are
$3,15,36,123,393,1302$.

This provides a useful resource estimate before constructing any
Hamiltonian: the Hilbert space grows as $d^L$, with $d>3$.
It also demonstrates why dimensions of categorical symmetry sectors
are not counts of low-energy primary fields.  These are dimensions
of entire finite lattice Hilbert spaces, including high-energy states.

\section{Four quantities that are sometimes all called an index}
\index{Jones index}

The original Haagerup subfactor has Jones index
\begin{equation}
 [M:N]=1+d=\frac{5+\sqrt{13}}2.
 \label{eq:original-Haagerup-index}
\end{equation}
Its relevant six-object even part is $\HH_2$, with a connected
algebra object $\one\oplus\rho$ of that dimension.  The
existence of this algebra structure is additional information beyond
the fusion rule.  In particular, one cannot infer the same algebra
structure in $\HH_3$ merely because its objects have the same
dimensions.

The CFT construction will instead use a chain of chiral extensions
\[
 U\subset B=V_{E_6\oplus A_2}\subset H=V_{E_8}.
\]
On the proper branch established by its construction inputs, the
extension dimensions are
\begin{center}
\begin{tabular}{c|c}
Inclusion or category & Relevant dimension\\\hline
Original Haagerup subfactor & $1+d$\\
$U\subset B$ & $s=1+d^2$\\
$B\subset H$ & $3$\\
$U\subset H$ & $3s$\\
$\HH_3$ & global dimension $3s$\\
$Z(\HH_3)$ & global dimension $(3s)^2$
\end{tabular}
\end{center}
The chiral extension entries will be justified later using quantum
dimensions and algebra objects.  The table is a notation guide,
not an assumption that the original subfactor inclusion literally
occurs inside this lattice VOA.  Multiplying the two extension
dimensions gives $3s$; replacing it by the original number $1+d$
would confuse two different constructions.

\section{Exercises}

\begin{exercise}\label{ex:H3-fusion}
Derive all four formulas in \eqref{eq:H3-all-products} from
\eqref{eq:H3-generating-ring}.  Verify the duality rule
\eqref{eq:H3-duals}, and give a two-line proof that the category
cannot be braided.

\par\smallskip\noindent{\footnotesize\normalfont
\hyperref[sol:ex:H3-fusion]{Solution on page~\pageref*{sol:ex:H3-fusion}.}}\par
\end{exercise}

\begin{exercise}\label{ex:H3-F}
Verify \eqref{eq:H3-F-pairs} directly from the radicals.  Use only
the three scalar identities \eqref{eq:H3-F-sum},
\eqref{eq:H3-F-pairs}, and \eqref{eq:H3-F-squares} to prove
orthogonality of the full matrix.  Explain why this calculation
does not prove every pentagon of $\HH_3$.

\par\smallskip\noindent{\footnotesize\normalfont
\hyperref[sol:ex:H3-F]{Solution on page~\pageref*{sol:ex:H3-F}.}}\par
\end{exercise}

\begin{exercise}\label{ex:H3-P}
Using the fusion ring alone, expand the square of the numerator in
\eqref{eq:H3-dimension-projector} and prove $P_+^2=P_+$.
Show that $Y_\rho P_+=dP_+$ without diagonalizing a matrix.

\par\smallskip\noindent{\footnotesize\normalfont
\hyperref[sol:ex:H3-P]{Solution on page~\pageref*{sol:ex:H3-P}.}}\par
\end{exercise}

\begin{exercise}\label{ex:H3-characteristic}
Compute the characteristic polynomial of
\eqref{eq:H3-path-adjacency}.  Use it to derive
\eqref{eq:H3-path-count} and obtain the length-six path count.

\par\smallskip\noindent{\footnotesize\normalfont
\hyperref[sol:ex:H3-characteristic]{Solution on page~\pageref*{sol:ex:H3-characteristic}.}}\par
\end{exercise}

\begin{exercise}\label{ex:H3-neutral}
Construct a unitary representation of the fusion algebra containing
three copies of the positive scalar representation, two copies of
the negative scalar representation, and four copies of the
two-dimensional representation.  What are the ranks of $E$, $P_+$,
$P_-$, and $P_{\rm ch}$?  How many states are neutral under
$\alpha$, and how many are fully invariant?

\par\smallskip\noindent{\footnotesize\normalfont
\hyperref[sol:ex:H3-neutral]{Solution on page~\pageref*{sol:ex:H3-neutral}.}}\par
\end{exercise}

\paragraph{Source note.}
The explicit associators follow the transparent presentation of
Huang and Lin.  The elementary fusion-algebra decomposition and
projector derivation also underlie the repository's symmetry-sector
checks.  Existing floating-point checks of the complete associator
are useful implementation checks; the hand calculation here proves
the displayed block identities exactly and claims no independent
reconstruction of the entire category.
